% page 291 of the facsimile (image p-290 of the scan)
% block 162.6 x 246.3 mm ; left margin 39.4 mm
\pgc{17.780mm}{0.000mm}{
\l{\cel{54}283}
\l{}
\l{}
\l{koercitiva. (\sou{elekt}r.)Feldo koercitiva : valoro de magneto-feldo,}
\l{\cel{3}necesa por retroduktar a zero la indukteso magnetala en ica}
\l{\cel{3}od ita punto di korpo feromagnetala, pos submisir ica a sat}
\l{\cel{3}multa cikli magnetala.}
\l{}
\l{kofio.\cel{2}Kapo-vesto mulierala, ek stofo, trikoturo, furo, e c.,}
\l{\cel{3}sen bordo. - EFS.}
\l{}
\l{kofro. Granda buxo ek ligno, ek metalo, rektangula, kun klozilo}
\l{\cel{3}klapa, maxim-multa-kaze konvexa, +kluzigebla per seruro.}
\l{\cel{3}- DEFIRS.}
\l{}
\l{koherar. (\sou{netrans}.)(\sou{aludante la elementi di ensemblo}).- Havar}
\l{\cel{3}sua parti strete unionita. - DEFIS.}
\l{}
\l{kohorto. (\sou{epoki antiq}ua) Trupo qua konsistis ye la dekimo di}
\l{\cel{3}legiono Romana. - DEFIRS.}
\l{}
\l{+\sou{koino}. Moneto-peci ek arjento, kupro, bronzo, nikelo, qui}
\l{\cel{3}valoras poke. (Kom ex.: Me kambiis moneto-peco 5-franka}
\l{\cel{3}po koino.)\cel{2}- E.}
\l{}
\l{koincidar\cel{3}(\sou{netrans}.) Interkontaktar precize, en omna punti.}
\l{\cel{3}- DEFIS.}
\l{}
\l{koitar. (\sou{netrans., kun}). (\sou{aludante homi})\cel{2}Unionar su sexuale,}
\l{\cel{3}por juo o genito. - DF.}
\l{}
\l{koko.\cel{2}Kombusteblo, reziduo de la min-karbono quan onu kalcinis}
\l{\cel{3}en vazo+kluza, por separar de olu la hido e la parti sulfa e}
\l{\cel{3}bitumina. - DEFIS.}
\l{}
\l{kokao.\cel{2}\sou{(bot.}) Arbusto aromatoza di Peru, di qua la folio konte-}
\l{\cel{3}nas elementi toniziva ed ecitigiva. - L.\cel{1}\sou{erythroxylum coca.}}
\l{\cel{3}DEFIRS.}
\l{}
\l{kokaino.\cel{2}(\sou{kemio)}.\cel{2}Alkaloido ek la kokao, quan onu uzas medicin}
\l{\cel{3}por anesteziar lokale. - DEFIRS.}
\l{}
\l{kokardo.\cel{2}Insigno, forme di mikra disko, qua +weresas an la}
\l{\cel{3}chapelo. - DEFIRS.}
\l{}
\l{koketa. Qua esforcas plezar a la personi di la altra sexuo.}
\l{\cel{3}- DEFIRS.}
\l{}
\l{kokleo. (\sou{zool.)}\cel{3}Molusko gasteropoda, reptera, kun konko helic-}
\l{\cel{3}atra. - L.\cel{1}\sou{cochlea.}\cel{1}- EFIS.}
\l{}
\l{koklusho. (patol.) Tusado konvulsatra, qua kaptas precipue la}
\l{\cel{3}infanti. - FR.}
\l{}
\l{kokono. Rnvelopilo quan facas, ek sua fili, la maxim multa}
\l{\cel{4}larvi (raupiI\cel{2}di la lepidoptari, ed en qua eventas lia last}
\l{\cel{4}metamorfoso. - DgFR.}
}
